\documentclass[10pt]{article}

\usepackage[utf8]{inputenc}

\usepackage[T1]{fontenc}

\usepackage{amsmath}

\usepackage{amsfonts}

\usepackage{amssymb}

\usepackage[version=4]{mhchem}

\usepackage{stmaryrd}

\usepackage{bbold}

\usepackage{mathrsfs}



\begin{document}

где $\alpha_{0}$ - высота траектории при $t=0$, a $\alpha^{\prime}$ - тангенс угла ее наклона к оси $t$ (это приближение оправдано, так как траектории Редже, как видно из рис. 4 , почти прямолинейны), то формулу (3) можно привести к виду (4), причем величина $B$ с увеличением энергии будет логарифмически расти:



\[

B(s)=B_{0}+2 \alpha^{\prime} \ln s,

\]



то есть рассеянные частицы с ростом энергии сосредоточиваются во все более узкой области передач импульса, так, как будто эффективный радиус $r$ сталкивающихся частиц растет:



\[

r^{2}=r_{0}^{2}+2 \alpha^{\prime} \ln s

\]



( $B_{0}$ и $r_{0}$ - величина наклона и радиус при $\left.s=1 \Gamma^{2} \mathrm{~B}^{2}\right)$.



Р. п. м. нашел широкое применение и в описании множественных процессов. В частности, в рамках этого метода естественно описываются такие явления, как скейлинг Фейнмана (см. Масштабная инвариантность), корреляция двух вторичных частиц. Одна из загадок физики элементарных частиц - наблюдаемая в эксперименте прямолинейность всех траекторий Редже и приблизительно одинаковые их наклоны.



Лит.: [1] Ш и рко в Д. В., «Успехи физич. наук», 1970, т. 102, в. 1; [2] Ко л л и н з П. Д. Б., С к в ай р с Э. Дж., Полюса Редже в физике частиц, пер. с англ., М., 1971. А. В. Ефремов, Д. В. Ширков. PEДКЕ полюсы в потенциальном рассеян и и - полюсы функции $S(\lambda, k)$, являющейся аналитическим продолжением чисел $S_{l}(k)$ - парциальных матриц рассеяния при фиксированной энергии $E=k^{2}-$ в плоскость комплексного углового момента $\lambda=l+\frac{1}{2}$. Впервые Р. п. в квантовой теории рассеяния исследованы Т. Редже [1] в рамках метода Ватсона преобразования. Кривые, описываемые Р. П. на плоскости комплексного углового момента при изменении энергии, называются трае ктор ия м и $\mathrm{Pe}-$ дже. При отрицательных энергиях $k^{2}<0$ Р. П. являются простыми и лежат на действительной оси $l$. При положительных $k^{2}>0$ траектории Редже, напротив, никогда не пересекают действительную ось.



Jum.: [1] Regge T., «Nuovo cimento», 1959, v. 14, p. 951-76; [2] Ne w t on R., The complex $j$-plane, N. Y.- Amst., 1964; [3] Нь юто н Р., Теория рассеяния волн и частиц, пер. с англ., М., 1969.



РЕДЖЕОН - см. Редже полюсов метод.



PEE - ШЛИДЕРА TEOPEМА - см. Алгебрачческий подход в Квантовой теории поля.



РЕЗЕРФОРДА ФОРМУЛА - формула для эффективного сечения рассеяния нерелятивистских заряженных точечных частиц, взаимодействующих по закону Кулона; получена Э. Резерфордом (Е. Rutherford, 1911). В системе центра инерции сталкивающихся частиц Р. Ф. имеет вид:



\[

\frac{d \sigma}{d \Omega}=\left(\frac{Z_{1} Z_{2} e^{2}}{2 m v^{2}}\right)^{2} \frac{1}{\sin ^{4}(\vartheta / 2)},

\]



где $d \sigma / d \Omega$ - сечение рассеяния в единичный телесный угол, $\vartheta$ - угол рассеяния, $m=m_{1} m_{2} /\left(m_{1}+m_{2}\right)$ - приведенная масса, $m_{1}, m_{2}-$ массы сталкивающихся частиц, $v-$ их относительная скорость, $Z_{1} e, Z_{2} e-$ электрич. заряды частиц $(e-$ элементарный электрич. заряд). Р.ф. справедлива как в классической, так и в квантовой теориях. Формула $\left({ }^{*}\right)$ была использована Э. Резерфордом при интерпретации опытов по рассеянию $\alpha$-частиц тонкими металлич. пластинками на большие углы $\left(\vartheta>90^{\circ}\right)$. В результате анализа опытных данных он пришел к выводу, что почти вся масса атома сконцентрирована в малом положительно заряженном ядре. Этим открытием были заложены основы современных представлений о строении атомов.



С. М. Биленький.



РЕЗОЛЬВЕНТА - см. Грина функция, Интегральное уравнение. PE30ЛЬBEHTA уравнен ия Л и у в лля - оператор, являющийся преобразованием Лапласа пропагатора уравнения Лиувилля. По определению, оператор Р. уравнения Лиувилля $\hat{R}(z), z \in \mathbb{C}$, задает временную эволюцию вектора распределения



\[

\vec{f}(t)=(2 \pi)^{-1} \int_{C} d z e^{-i z t} \hat{\mathbb{R}}(z) \vec{f}(0)

\]



где $C-$ контур в комплексной плоскости $z$, образованный прямой, проходящей параллельно действительной оси выше всех сингулярностей функции $\hat{\Re}(\mathrm{z})$. Р. уравнения Лиувилля и пропагатор $\hat{U}(t)$ связаны соотношением



\[

\hat{U}(t)=(2 \pi)^{-1} \int_{C} d z e^{-i z t} \hat{\Re}(z) .

\]



Р. уравнения Лиувилля формально можно определить в виде $\hat{\mathscr{A}}(z)=(-\hat{\mathscr{L}}-i z)^{-1}$, где $\hat{\mathscr{\mathscr { }}}-$ оператор Лиувилля.



Лит.: [1] Балес ку Р., Равновесная и неравновесная статистическая механика, пер. с англ., т. 2, М., 1978; [2] Prigogine I., Ge orge C., He ni n F., «Physica», 1969, v. 45, p. 418.



Д. П. Санкович.



PEЗОНАНС - явление увеличения амплитуды вынужденных колебаний при приближении частоты внешнего воздействия к одной из частот собственных колебаний динамической системы. Явление Р. имеет наиболее простой характер в линейной динамич. системе. Дифференциальное уравнение движения линейной системы с одной степенью свободы в среде с вязким трением при гармонич. воздействии имеет вид:



\[

a \ddot{q}+b \dot{q}+c q=H \sin (p t+\delta),

\]



где $q$ - обобщенная координата; $a, b, c$ - постоянные параметры, характеризуюшие систему; $\boldsymbol{H}, \boldsymbol{p}, \boldsymbol{\delta}$ - соответственно амплитуда, частота, начальная фаза внешнего воздействия. Установившиеся вынужденные колебания происходят по гармонич. закону с частотой $p$ и амплитудой



\[

D=\frac{H}{a \sqrt{\left(k^{2}-p^{2}\right)^{2}+\frac{b^{2}}{a^{2}} p^{2}}},

\]



где $k=\sqrt{c / a}-$ частота собственных колебаний при отсутствии рассеивания энергии $(b=0)$. Амплитуда $D$ имеет максимальное значение при



\[

\frac{p}{k}=\sqrt{1-\frac{b^{2}}{2 a c}}

\]



и при малом рассеивании энергии близка к ее значению при $p=k$. Принято называть ре з о на н с о м тот случай, когда $p=k$. Если $b=0$, то при $p=k$ амплитуда вынужденных колебаний возрастает пропорционально времени.



Если линейная система имеет $n$ степеней свободы, то P. наступает при совпадении частоты внешней силы с одной из собственных частот системы. При негармонич. воздействии Р. может иметь место лишь при совпадении частот его гармонич. спектра с частотами собственных колебаний.



См. также Малые знаменатели, Обмотка тора, Осциллятор, Резонансное соотношение.



Лит.: [1] С т ре л к о в С. П., Введение в теорию колебаний, М.$\begin{array}{ll}\text { Л., 1951. } & \text { Н. В. Бутенин. }\end{array}$



PEЗОНАНСНАЯ ЗОНА - см. КАМ-теория.



PE3OHAHCHOE СOOTHOШEHИE, р е з о н а н с,- 1) для системы с быстро вращающимися фазами - равенство, выражающее целочисленную линейную зависимость частот: $(\boldsymbol{k}, \boldsymbol{\omega})=0$, где $\boldsymbol{k}=\left(k_{1}, \ldots, k_{\boldsymbol{m}}\right)$ - целочисленный ненулевой вектор с несократимыми компонентами, $\omega=\omega(I)=$ $=\left(\omega_{1}(I), \ldots, \omega_{m}(I)\right)-$ вектор частот невозмущенного движения, $(\boldsymbol{k}, \omega) \equiv k_{1} \omega_{1}+\ldots+k_{m} \omega_{m}$. Число $|\boldsymbol{k}|=\left|k_{1}\right|+\ldots+\left|k_{m}\right|$ называется порядком Р.с. (или порядком резона н с а). Гармоники



\[

\sin _{\cos }[i l(k, \varphi)], \quad l \in \mathbb{Z} \backslash\{0\}, \varphi \operatorname{modd} 2 \pi \in T^{m}

\]



называются резонансным и гармоникам и данного Р. с. Поверхность в пространстве $\mathbb{R}^{n}$ медленных переменных





\end{document}